\section{Related work}
\label{sec:related}

We situate \pqlang{} among quantum programming languages and frameworks,
then among algorithm libraries, and finally record the feature comparison of
Table~\ref{tab:qpl-comparison} and the coverage comparison of
Table~\ref{tab:ode-coverage}.  The discussion is organized by the
requirements R1--R8 of \secref{sec:qsc-requirements}; general surveys of the
field are available~\cite{selinger2004towards,heim2020quantum,garhwal2021systematic}.

\subsection{Quantum programming languages}
\label{sec:related-qpl}

\paragraph{Functional and typed language lines.}
Quipper~\cite{green2013quipper} pioneered scalable, circuit-generating
functional programming: a Haskell host generates gate-level circuits, with
``boxing'' providing a form of abstraction over generated subcircuits.
Scaffold~\cite{javadiabhari2015scaffcc} embedded modules and classical
control in a C-like surface, compiled by ScaffCC with early
resource-estimation passes.  Notably, Quipper later served as the vehicle
for landmark resource analyses of quantum linear systems---although as
external studies over generated circuits, not as a language
mechanism~\cite{scherer2017resource,fillion2017dirac}.  QWIRE and the Proto-Quipper line developed linear and dependent type
disciplines for safety guarantees~\cite{paykin2017qwire,fu2023protoquipper}.
Silq~\cite{bichsel2020silq} introduced safe automatic uncomputation through
its type system---a genuine language-level answer to the invertibility half
of our R4.  None of these lines, however, provides deferred-oracle
semantics with paradigm contracts and late linking (R1/R2): abstractions
expand to circuits, and the open state of a program is not a shareable
artifact.  Their IRs target gate-level compilation rather than
module-preserving program storage (R3/R5), and none ships a
scientific-computing algorithm stack.

\paragraph{Industrial SDKs and embedded frameworks.}
ProjectQ~\cite{steiger2018projectq} established the Python-embedded pattern,
shipped one of the earliest resource estimators, and---through its
extensible compiler-engine chain---explicitly supports \emph{emulating}
quantum programs at a high level of abstraction, mimicking large oracles
instead of compiling them to gates: a genuine form of deferred oracle at
execution time, though one that lives in the engine stack rather than in a
serializable open program.  Quil and pyQuil~\cite{smith2016quil} and
Cirq~\cite{cirqdevelopers2020cirq} provide circuit DSLs centered on NISQ execution.
Qiskit~\cite{javadiabhari2024qiskit} grew a large circuit library and textbook algorithm
implementations; demonstration-level HHL shipped historically in its
algorithms module and has since been retired with it, and no QSVT or
block-encoding abstraction exists in the library today.  Its objects are
concrete circuits, its transpiler flattens structure, and alternative oracle
realizations are separate programs rather than bindings of one open slot.
OpenQASM~3~\cite{cross2022openqasm3} standardizes a textual interchange
language whose subroutines and loops preserve some structure, and whose
``opaque'' declarations gesture toward R1---but at gate granularity,
without paradigm contracts, capabilities, or a linking discipline.
Closest to our primitive layer among embedded frameworks, Qrisp~\cite{seidel2023qrisp}
automates uncomputation and resource-friendly scheduling inside Python and
exposes a \texttt{BlockEncoding} class---block-encodings as programming
abstractions with a {NumPy}-like interface---which we regard as independent
evidence that scientific-computing primitives belong at the language level.
Qrisp does not, however, treat the open program as an artifact: there is no
paradigm-contracted oracle declaration with late, partial, repeatable
linking, no module-preserving serialized IR with promise provenance, and no
ODE/PDE solver families.  t\textbar ket\textgreater~\cite{sivarajah2021tket} and
staq~\cite{amy2020staq} are compiler and synthesis toolchains,
complementary to our concerns: they operate \emph{below} the register/module
level our RIR preserves.

\paragraph{Q\#.}
Q\#~\cite{svore2018qsharp} is the closest prior art at the language level.
Its operations with \texttt{Adjoint}/\texttt{Controlled} \emph{functors}
anticipate our capability discipline (R4): specialization availability is
declared and tracked.  Passing operations as parameters---the closest
mainstream analogue to our oracle slots---lets algorithms be written against
any conforming implementation, so Q\# partially meets R1 at \emph{call}
time.  Its standard libraries include fixed-point numerics and estimation
primitives, and the surrounding Azure tooling ships a serious resource
estimator.  What Q\# does not provide for our workload is
the \emph{open program}: oracles are operations with concrete bodies;
``partial application'' is host-level; there is no serialized artifact in
which an algorithm carries its unimplemented input models with paradigm,
$\alpha$, capabilities, and provenance (R1/R2/R6); and its libraries do not
include block-encoding algebra, a QSVT/QSP engine, or any ODE/PDE solver
families (Table~\ref{tab:ode-coverage}).  The distinction is one of
\emph{when} semantics attach: Q\# resolves capabilities at compile time of
concrete operations, whereas \pqlang{} keeps the \emph{link} itself as a
first-class, repeatable, partial operation on stored programs.

\paragraph{Research codebases for scientific algorithms.}
Individual algorithm implementations---LCHS, Schr\"odingerization, Carleman
pipelines, QLSS prototypes---exist as research code accompanying
publications~\cite{an2026lchs,jin2024schrodingerization,liu2021pnas,
liu2023carleman,an2026lchs}.  Each typically fixes one method, one input model, and one
small simulator; none offers shared input-model objects, interchangeable
kernels, or a common IR.  The absence of a unifying language layer is
precisely the gap this paper documents: the algorithms compose
mathematically (all consume block-encodings; all produce state oracles) but
not computationally, absent exactly the contracts of
\secref{sec:oracles}.

\begin{table}[t]
  \centering
  \caption{Language and framework comparison against the requirements of
  \secref{sec:qsc-requirements}.  \ok~= provided; \half~= partial (see
  text); \noent~= absent.  Claims are referenced to the systems'
  publications and documentation.}
  \label{tab:qpl-comparison}
  \small
  \setlength{\tabcolsep}{4pt}\begin{tabular}{@{}p{2.4cm}cccccc@{}}
    \toprule
    & \textbf{Quipper} & \textbf{Scaffold} & \textbf{Q\#} & \textbf{ProjectQ} & \textbf{Qiskit} & \textbf{\pqlang} \\
    \midrule
    R1 deferred oracles & \noent$^{a}$ & \noent & \half$^{b}$ & \half$^{c}$ & \noent & \ok \\
    R2 realization independence & \noent & \noent & \half & \noent$^{c}$ & \half & \ok \\
    R3 structure-preserving IR & \half & \half & \half & \noent & \noent & \ok \\
    R4 adjoint/control tracked & \half & \noent & \ok & \noent & \noent & \ok \\
    R5 symbolic scale ($2^{40}$+) & \half & \half & \half & \noent & \noent & \ok \\
    R6 promise provenance & \noent & \noent & \noent & \noent & \noent & \ok \\
    R7 readout separation & \half & \half & \noent & \noent & \noent & \ok \\
    R8 library extensibility & \half & \half & \half & \half & \half & \ok \\
    \midrule
    Scientific primitives$^{d}$ & \noent & \noent & \half & \noent & \half & \ok \\
    ODE/PDE solver families & \noent & \noent & \noent & \noent & \noent & \ok \\
    Resource estimation & \half$^{e}$ & \ok & \ok & \ok & \half & \noent$^{f}$ \\
    \bottomrule
    \multicolumn{7}{@{}p{14.3cm}@{}}{\scriptsize $^a$Boxes abstract over generated subcircuits, but there is no open-program state.  $^b$Operations are passed as parameters and remain abstract to the algorithm, but bodies are concrete and no link-time openness or paradigm contracts exist.  $^c$The extensible engine chain can emulate oracles at execution time, but not in a serializable program artifact.  $^d$Block-encoding algebra,
    QSVT/QSP engine, sparse-access and QRAM input models, reversible
    fixed-point arithmetic with a math-function compiler.
    $^f$Scoped and reserved (\secref{sec:reserved}); the IR preserves the
    structure an estimator needs.}
  \end{tabular}
\end{table}

\begin{table}[t]
  \centering
  \caption{Algorithm and solver coverage: \pqlang{} against the major
  frameworks' standard libraries (research-prototype implementations
  published outside a framework's libraries do not count).  \ok~= provided
  as a composable, contract-checked library component; \half~=
  demonstration-level implementation; \noent~= absent.}
  \label{tab:ode-coverage}
  \small
  \setlength{\tabcolsep}{4pt}\begin{tabular}{@{}p{5.4cm}ccccc@{}}
    \toprule
    \tabhead{Capability} & \textbf{Q\#} & \textbf{ProjectQ} & \textbf{Qiskit} & \textbf{Cirq} & \textbf{\pqlang} \\
    \midrule
    Query/search/amplification & \ok & \ok & \ok & \ok & \ok \\
    Phase/amplitude estimation & \ok & \ok & \ok & \half & \ok \\
    Hamiltonian simulation (Trotter) & \ok & \half & \ok & \half & \ok \\
    Hamiltonian simulation (LCU/Taylor, QSVT) & \noent & \noent & \half & \noent & \ok \\
    Block-encoding algebra & \noent & \noent & \half & \noent & \ok \\
    QSP phase synthesis (self-verifying) & \noent & \noent & \noent & \noent & \ok \\
    Quantum linear systems (QLSS) & \noent & \noent & \half$^{a}$ & \noent & \ok \\
    QRAM/QROM input models (typed resource) & \noent & \noent & \half & \noent & \ok \\
    Reversible fixed-point arithmetic & \ok & \noent & \noent & \noent & \ok \\
    Math-function compiler (Python $\rightarrow$ reversible) & \noent & \noent & \noent & \noent & \ok \\
    ODE: LCHS & \noent & \noent & \noent & \noent & \ok \\
    ODE: Schr\"odingerization & \noent & \noent & \noent & \noent & \ok \\
    ODE: contour decomposition (CBMD) & \noent & \noent & \noent & \noent & \ok \\
    ODE: Carleman (nonlinear) & \noent & \noent & \noent & \noent & \ok \\
    ODE: history systems via QLSS & \noent & \noent & \noent & \noent & \ok \\
    PDE: heat / Fokker--Planck input models & \noent & \noent & \noent & \noent & \ok \\
    PDE: symbolic nonlinear PDEs (QHAM) & \noent & \noent & \noent & \noent & \ok \\
    PDE: Euler equations (QFVM/Roe) & \noent & \noent & \noent & \noent & \ok \\
    \bottomrule
    \multicolumn{6}{@{}p{14.0cm}@{}}{\scriptsize $^a$A demonstration-level HHL circuit shipped historically in Qiskit's (since-retired) algorithms module; it is not part of the current library.}
  \end{tabular}
\end{table}

\subsection{Positioning}
\label{sec:positioning}

The comparison is not a claim that prior systems are poorly built; they are
built for a different workload.  Circuit SDKs optimize for NISQ execution
and hardware co-design; typed language lines optimize for semantic safety
of full quantum computation; Q\# optimizes for a vertically integrated
development kit.  \pqlang{} optimizes for the fault-tolerant
scientific-computing workload, where the unit of progress is a
\emph{mathematical algorithm against declared access models}, the unit of
sharing is a \emph{program artifact with open slots and recorded promises},
and the unit of validation is \emph{graduated evidence} rather than a single
simulator's output.  To our knowledge, no previously published system
combines these three units; Tables~\ref{tab:qpl-comparison}
and~\ref{tab:ode-coverage} make the comparison concrete, and the honest
counterpoint is equally explicit: resource estimation---strong in Scaffold,
Q\# tooling, and ProjectQ---is the deliberately reserved axis of
\secref{sec:reserved}.
